\documentclass[11pt,a4paper]{article}

% ---------------------------------------------------------
% PREAMBLE (stable, warning-free, Overleaf-safe)
% ---------------------------------------------------------
\usepackage[utf8]{inputenc}
\usepackage[T1]{fontenc}
\usepackage{lmodern}
\usepackage{amsmath, amssymb, amsfonts}
\usepackage{bm}
\usepackage{geometry}
\usepackage{graphicx}
\usepackage{hyperref}
\usepackage{cite}
\usepackage{authblk}
\usepackage{physics}
\usepackage{mathtools}

\geometry{margin=2.4cm}

\title{\textbf{Companion LXXXI: Diffusion–Optimal Transport Unification and Neural Transport Maps for Persistence Diagrams in the Relaxation-Driven Cyclic Cosmology}}
\author{Michael Lehmann \\ \texttt{mi.lehmann@gmx.de}}
\date{April 2026}

\begin{document}
\maketitle

\begin{abstract}
Diffusion models and optimal transport (OT) provide two complementary perspectives 
on probability distributions: diffusion models describe generative stochastic 
dynamics, while OT characterizes deterministic transport maps. Recent advances 
have shown that these two viewpoints can be unified through stochastic control, 
Schrödinger bridges, and score-based transport geometry. Building on the neural 
Schrödinger flows developed in Companion LXXX, this work introduces a unified 
diffusion–OT framework for persistence diagrams in the Relaxation-Driven Cyclic 
Cosmology (RDCC). The relaxation mechanism modifies the scale-dependent evolution 
of the gravitational potential, inducing characteristic changes in both the 
deterministic OT map and the stochastic diffusion flow. The resulting neural 
transport maps provide a powerful tool for modeling RDCC signatures, enabling 
likelihood-free inference, generative modeling, and transport-based parameter 
constraints from weak-lensing topology.
\end{abstract}
% ---------------------------------------------------------
\section{Introduction}

Persistence diagrams encode the topological structure of weak-lensing convergence 
fields. Previous Companions introduced Wasserstein geometry, entropic optimal 
transport, Schrödinger bridges, score-based diffusion models, and neural 
Schrödinger flows as tools for analyzing the Relaxation-Driven Cyclic Cosmology 
(RDCC).

Diffusion models provide stochastic generative dynamics, while optimal transport 
provides deterministic transport maps. Unifying these perspectives yields a 
powerful framework in which:

\begin{itemize}
    \item diffusion models approximate OT maps in the zero-noise limit,
    \item OT maps arise as the mean flow of Schrödinger bridges,
    \item score fields encode the gradient of the log-density,
    \item neural SDEs learn both stochastic and deterministic transport geometry.
\end{itemize}

In RDCC, the relaxation mechanism modifies both the deterministic and stochastic 
components of transport, producing characteristic signatures in the unified 
diffusion–OT geometry.
% ---------------------------------------------------------
\section{Optimal Transport Maps for Persistence Diagrams}

Let $\rho_{0}$ and $\rho_{1}$ be distributions of persistence diagrams. The OT map 
$T$ solves

\begin{equation}
    T = \arg\min_{T}
    \int \| D - T(D) \|^{2} \, d\rho_{0}(D),
\end{equation}

subject to $T_{\#} \rho_{0} = \rho_{1}$.

For persistence diagrams, $T$ acts on birth-death coordinates and encodes the 
deterministic geometric deformation between topological ensembles.
% ---------------------------------------------------------
\section{Diffusion Models as Stochastic Transport Maps}

A diffusion model defines the reverse-time SDE

\begin{equation}
    dD_{t}
    = - \beta(t) s(D_{t}, t) \, dt
      + \sqrt{\beta(t)} \, d\bar{B}_{t}.
\end{equation}

The mean flow of this SDE approximates the OT map:

\begin{equation}
    \mathbb{E}[D_{t}] \approx T(D_{0}).
\end{equation}

Thus diffusion models provide a stochastic relaxation of OT maps.
% ---------------------------------------------------------
\section{Unified Diffusion–OT Geometry}

The unified framework is governed by the controlled SDE

\begin{equation}
    dD_{t}
    = \left[
        - \varepsilon(t) s(D_{t}, t)
        + u(D_{t}, t)
      \right] dt
      + \sqrt{2 \varepsilon(t)} \, dB_{t},
\end{equation}

where

\begin{itemize}
    \item $s(D_{t}, t)$ is the score field,
    \item $u(D_{t}, t)$ is a deterministic OT-like control,
    \item $\varepsilon(t)$ interpolates between OT ($\varepsilon \to 0$) and diffusion.
\end{itemize}

This yields a continuum between:



\[
\text{deterministic OT} \quad \leftrightarrow \quad \text{stochastic diffusion}.
\]


% ---------------------------------------------------------
\section{RDCC Modification of Unified Transport Geometry}

In RDCC, the convergence evolves as

\begin{equation}
    \kappa_{\mathrm{RDCC}}(\ell)
    = \kappa(\ell)
      \left[
        1 - \chi_{\kappa}
        \left(\frac{\ell}{\ell_{\mathrm{rel}}}\right)^{\alpha}
      \right].
\end{equation}

This modifies:

\begin{itemize}
    \item the score field $s(D,t)$,
    \item the deterministic OT control $u(D,t)$,
    \item the diffusion schedule $\varepsilon(t)$,
    \item the curvature of transport trajectories,
    \item the barycentric flow across redshift.
\end{itemize}

RDCC induces:

\begin{itemize}
    \item enhanced deterministic transport at large scales,
    \item suppressed stochasticity at small scales,
    \item a characteristic turnover near $\ell_{\mathrm{rel}}$,
    \item scale-dependent deformation of transport maps.
\end{itemize}
% ---------------------------------------------------------
\section{Conclusion}

Diffusion models and optimal transport provide complementary perspectives on the 
geometry of probability distributions. Neural Schrödinger flows unify these 
perspectives into a single framework for learning stochastic and deterministic 
transport processes. In the Relaxation-Driven Cyclic Cosmology, the relaxation 
mechanism modifies both components of transport in a scale-dependent manner, 
producing distinctive signatures in weak-lensing topology. The present Companion 
establishes the foundation for transport-based RDCC parameter inference, 
generative modeling, and unified geometric analysis of topological data.

\bibliographystyle{apsrev4-2}
\nocite{*}
\bibliography{references}


\end{document}
